%%% LaTeX-Vorlage Version 1.4 %%%
\documentclass[a4paper,fontsize=11pt,oneside,parskip=half,headings=normal]{scrreprt} 
% \usepackage{showframe} % nur für Kontrolle der Ränder 

%%% Präambel einbinden %%%
%%% Präambel %%%
% hier sollten keine Änderungen erforderlich sein
%
\usepackage[utf8]{inputenc}   % Zeichencodierung UTF-8 für Eingabe-Dateien
\usepackage[T1]{fontenc}      % Darstellung von Umlauten im PDF

\usepackage{svg}
\usepackage{tabularx}
\usepackage[bottom]{footmisc}
\usepackage{listings}         % für Einbindung von Code-Listings
\lstset{numbers=left,numberstyle=\tiny,numbersep=5pt,texcl=true}
\lstset{literate=             % erlaubt Sonderzeichen in Code-Listings 
{Ö}{{\"O}}1
{Ä}{{\"A}}1
{Ü}{{\"U}}1
{ß}{{\ss}}2
{ü}{{\"u}}1
{ä}{{\"a}}1
{ö}{{\"o}}1
{€}{{\euro}}1
}

\usepackage[
  inner=35mm,outer=15mm,top=25mm,
  bottom=20mm,foot=12mm,includefoot
]{geometry}                 % Einstellungen für Ränder

\usepackage[english]{babel} % Spracheinstellungen Englisch
\usepackage[babel,english=british]{csquotes} % deutsche Anf.zeichen
\usepackage{enumerate}      % anpassbare Nummerier./Aufz.
\usepackage{graphicx}       % Einbinden von Grafiken
\usepackage[onehalfspacing]{setspace} % anderthalbzeilig

\usepackage{blindtext}      % Textgenerierung für Testzwecke
\usepackage{color}          % Verwendung von Farbe 

\usepackage{acronym}        % für ein Abkürzungsverzeichnis

\usepackage[                % Biblatex
  backend=biber,
  bibstyle=_dhbw_authoryear,
  citestyle=authoryear,
  uniquename=true, useprefix=true,
  bibencoding=utf8]{biblatex}
%kein Punkt am Ende bei \footcite
%http://www.golatex.de/footcite-ohne-punkt-am-schluss-t4865.html
\renewcommand{\bibfootnotewrapper}[1]{\bibsentence#1}


%Reihenfolge der Autorennamen
%   
% http://golatex.de/viewtopic,p,80448.html#80448
% Argumente: siehe http://texwelt.de/blog/modifizieren-eines-biblatex-stils/
\DeclareNameFormat{sortname}{% Bibliographie
  \ifnum\value{uniquename}=0 % Normalfall
    \ifuseprefix%
      {%
         \usebibmacro{name:family-given}
           {\namepartfamily}
           {\namepartgiveni}
           {\namepartprefix}
           {\namepartsuffixi}%
       }
      {%
         \usebibmacro{name:family-given}
           {\namepartfamily}
           {\namepartgiveni}
           {\namepartprefixi}
           {\namepartsuffixi}%
       }%
  \fi
  \ifnum\value{uniquename}=1% falls nicht eindeutig, abgek. Vorname 
      {%
         \usebibmacro{name:family-given}
           {\namepartfamily}
           {\namepartgiveni}
           {\namepartprefix}
           {\namepartsuffix}%
       }%
  \fi
  \ifnum\value{uniquename}=2% falls nicht eindeutig, ganzer Vorname 
      {%
         \usebibmacro{name:family-given}
           {\namepartfamily}
           {\namepartgiven}
           {\namepartprefix}
           {\namepartsuffix}%
       }%
  \fi   
  \usebibmacro{name:andothers}}

\DeclareNameFormat{labelname}{% für Zitate
  \ifnum\value{uniquename}=0 % Normalfall
    \ifuseprefix%
      {%
         \usebibmacro{name:family-given}
           {\namepartfamily}
           {\empty}
           {\namepartprefix}
           {\namepartsuffixi}%
       }
      {%
         \usebibmacro{name:family-given}
           {\namepartfamily}
           {\empty}
           {\namepartprefixi}
           {\namepartsuffixi}%
       }%
  \fi
  \ifnum\value{uniquename}=1% falls nicht eindeutig, abgek. Vorname 
      {%
         \usebibmacro{name:family-given}
           {\namepartfamily}
           {\namepartgiveni}
           {\namepartprefix}
           {\namepartsuffix}%
       }%
  \fi
  \ifnum\value{uniquename}=2% falls nicht eindeutig, ganzer Vorname 
      {%
         \usebibmacro{name:family-given}
           {\namepartfamily}
           {\namepartgiven}
           {\namepartprefix}
           {\namepartsuffix}%
       }%
  \fi   
  \usebibmacro{name:andothers}}
      
  
\DeclareFieldFormat{extrayear}{% = the 'a' in 'Jones 1995a'
  \iffieldnums{labelyear}
    {\mknumalph{#1}}
    {\mknumalph{#1}}}        

\renewcommand*{\multinamedelim}{\addslash}
\renewcommand*{\finalnamedelim}{\addslash}
\renewcommand*{\multilistdelim}{\addslash}
\renewcommand*{\finallistdelim}{\addslash}

\renewcommand{\nameyeardelim}{~}

% Literaturverzeichnis: Doppelpunkt zwischen Name (Jahr): Rest 
% http://de.comp.text.tex.narkive.com/Tn1HUIXB/biblatex-authoryear-und-doppelpunkt
\renewcommand{\labelnamepunct}{\addcolon\addspace}

% Im Literaturverzeichnis: Autor (Jahr) fett
\renewbibmacro*{author}{%
  \ifboolexpr{%
    test \ifuseauthor%
    and
    not test {\ifnameundef{author}}
  }
    {\usebibmacro{bbx:dashcheck}
       {\bibnamedash}
       {\usebibmacro{bbx:savehash}%
        \textbf{\printnames{author}}%
        \iffieldundef{authortype}
          {\setunit{\addspace}}
          {\setunit{\addcomma\space}}}%
     \iffieldundef{authortype}
       {}
       {\usebibmacro{authorstrg}%
        \setunit{\addspace}}}%
    {\global\undef\bbx@lasthash
     \usebibmacro{labeltitle}%
     \setunit*{\addspace}}%
  \textbf{\usebibmacro{date+extrayear}}}

% Sonderfall: Quelle ohne Autor, aber mit Herausgeber
% Name des Herausgebers wird fett gedruckt
\renewbibmacro*{bbx:editor}[1]{%
  \ifboolexpr{%
    test \ifuseeditor%
    and
    not test {\ifnameundef{editor}}
  }
    {\usebibmacro{bbx:dashcheck}
       {\bibnamedash}
       {\textbf{\printnames{editor}}%
        \setunit{\addcomma\space}%
        \usebibmacro{bbx:savehash}}%
     \usebibmacro{#1}%
     \clearname{editor}%
     \setunit{\addspace}}%
    {\global\undef\bbx@lasthash
     \usebibmacro{labeltitle}%
     \setunit*{\addspace}}%
  \textbf{\usebibmacro{date+extrayear}}}

% Anpassungen für deutsche Sprache
\DefineBibliographyStrings{ngerman}{%
	nodate = {{o.J.}},
	urlseen = {{Abruf:}},
	ibidem = {{ebenda}}
}

% Benennung der Verzeichnisse
\defbibheading{lit}{\section*{List of literature}}
\defbibheading{www}{\section*{List of Internet and intranet sources}}
\defbibheading{gespraeche}{\section*{List of conversations}}

% Anpassungen für die englische Sprache
\DefineBibliographyStrings{english}{
    nodate = {{w.y.}},
    urlseen = {{retrieval:}}
} % mit DHBW-spezifischen Einstellungen

\usepackage{hyperref}       % URL-Formatierung, klickbare Verweise

\usepackage{tocloft}        % für Verzeichnis der Anhänge

\newcounter{anhcnt}
\setcounter{anhcnt}{0}
\newlistof{anhang}{app}{}

\newcommand{\anhang}[1]{%
  \refstepcounter{anhcnt}
  \setcounter{anhteilcnt}{0}
  \section*{Appendix \theanhcnt: #1}
  \addcontentsline{app}{section}{\protect\numberline{A.\theanhcnt}#1}\par
}

\newcounter{anhteilcnt}
\setcounter{anhteilcnt}{0}

\newcommand{\anhangteil}[1]{%
	\refstepcounter{anhteilcnt}
	\subsection*{Appendix~\arabic{anhcnt}/\arabic{anhteilcnt}: #1}
	\addcontentsline{app}{subsection}{\protect\numberline{A.\theanhcnt/\arabic{anhteilcnt}}#1}\par
}

\renewcommand{\theanhteilcnt}{Appendix \theanhcnt/\arabic{anhteilcnt}}

\usepackage{chngcntr}                % fortlaufende Zähler für Fußnoten, Abbildungen und Tabellen
\counterwithout{figure}{chapter}
\counterwithout{table}{chapter}
\counterwithout{footnote}{chapter}

\usepackage[automark]{scrlayer-scrpage} 
%% Definitionen für Kopf- und Fußzeile auf normalen Seiten
\defpagestyle{kopfzeile}
{% Kopfdefinition
  (\textwidth,0pt)    % Länge der oberen Linie,Dicke der oberen Linie       
  {} % Definition für linke Seiten im doppelseitigen Layout
  {} % Definition für rechte Seiten im doppelseitigen Layout      
  {  % Definition für Seiten im einseitigen Layout
	\makebox[0pt][l]{\rightmark}% 
	\makebox[\linewidth]{}% 
  }        
  (\textwidth, 0.4pt) % Untere Linienlänge, Untere Liniendicke
}
{% Fußdefinition
  (\textwidth,0pt)    % Obere Linienlänge, Obere Liniendicke
  {} % Definition für linke Seiten im doppelseitigen Layout
  {} % Definition für rechte Seiten im doppelseitigen Layout
  {  % Definition für Seiten im einseitigen Layout
    \makebox[\linewidth]{}%
    \makebox[0pt][r]{\pagemark}%
  }
  (\textwidth, 0pt)   % Länge der unteren Linie,Dicke der unteren Linie
}

%% Definitionen für Kopf- und Fußzeile auf ersten Seiten eines Kapitels
\defpagestyle{kapitelkopfzeile}
{% Kopfdefinition
  (\textwidth,0pt)    % Länge der oberen Linie,Dicke der oberen Linie       
  {} % Definition für linke Seiten im doppelseitigen Layout
  {} % Definition für rechte Seiten im doppelseitigen Layout      
  {}  % Definition für Seiten im einseitigen Layout
  (\textwidth, 0pt) % Untere Linienlänge, Untere Liniendicke
}
{% Fußdefinition
  (\textwidth,0pt)    % Obere Linienlänge, Obere Liniendicke
  {} % Definition für linke Seiten im doppelseitigen Layout
  {} % Definition für rechte Seiten im doppelseitigen Layout
  {  % Definition für Seiten im einseitigen Layout
    \makebox[\linewidth]{}%
    \makebox[0pt][r]{\pagemark}%
  }
  (\textwidth, 0pt)   % Länge der unteren Linie,Dicke der unteren Linie
}

%% Definitionen für Kopf- und Fußzeile im Anhang und bei Quellenverzeichnisse
\newcommand{\spezialkopfzeileBezeichnung}{}
\defpagestyle{spezialkopfzeile}
{% Kopfdefinition
  (\textwidth,0pt)    % Länge der oberen Linie,Dicke der oberen Linie       
  {} % Definition für linke Seiten im doppelseitigen Layout
  {} % Definition für rechte Seiten im doppelseitigen Layout      
  {  % Definition für Seiten im einseitigen Layout
	\makebox[0pt][l]{\spezialkopfzeileBezeichnung}% 
	\makebox[\linewidth]{}% 
  }        
  (\textwidth, 0.4pt) % Untere Linienlänge, Untere Liniendicke
}
{% Fußdefinition
  (\textwidth,0pt)    % Obere Linienlänge, Obere Liniendicke
  {} % Definition für linke Seiten im doppelseitigen Layout
  {} % Definition für rechte Seiten im doppelseitigen Layout
  {  % Definition für Seiten im einseitigen Layout
    \makebox[\linewidth]{}%
    \makebox[0pt][r]{\pagemark}%
  }
  (\textwidth, 0pt)   % Länge der unteren Linie,Dicke der unteren Linie
}
            
\newcommand\spezialkopfzeile[1]{%
  \renewcommand\spezialkopfzeileBezeichnung{#1}
  \pagestyle{spezialkopfzeile}
}                 
                  
% Standard-Pagestyle auswählen
\pagestyle{kopfzeile}

% keine Kopfzeile anzeigen auf Seiten, auf denen ein 
% Kapitel beginnt oder das Inhalts-/Abbildungs-/Tabellenverzeichnis steht 
\renewcommand{\chapterpagestyle}{kapitelkopfzeile}
\tocloftpagestyle{kapitelkopfzeile}

		 % für schöne Kopfzeilen 

\usepackage{textcomp}			 % erlaubt EUR-Zeichen in Eingabedatei
\usepackage{eurosym}             % offizielles EUR-Symbol in Ausgabe
\renewcommand{\texteuro}{\euro}  % ACHTUNG: nach hyperref aufrufen!

\usepackage{scrhack}             % stellt Kompatibilität zw. KOMA-Script
                                 % (scrreprt) und anderen Paketen her
                                 
% Anpassung der Abstände bei Kapitelüberschriften
% (betrifft auch Inhalts-, Abbildungs- und Tabellenverzeichnis)
\renewcommand*\chapterheadstartvskip{\vspace*{-\topskip}}
\newcommand{\myBeforeTitleSkip}{1mm}
\newcommand{\myAfterTitleSkip}{10mm}
\setlength\cftbeforetoctitleskip{\myBeforeTitleSkip}
\setlength\cftbeforeloftitleskip{\myBeforeTitleSkip}
\setlength\cftbeforelottitleskip{\myBeforeTitleSkip}

\setlength\cftaftertoctitleskip{\myAfterTitleSkip}
\setlength\cftafterloftitleskip{\myAfterTitleSkip}
\setlength\cftafterlottitleskip{\myAfterTitleSkip}                                                            
%%% Ende der Präambel %%%

%%% Name Ihrer Literatur-Datenbank (ggf. anpassen) %%%
\bibliography{includes/literatur-datenbank.bib}

\begin{document}
%%% Deckblatt einbinden %%% 
% Anpassungen nötig (Name, Titel etc.)
% HIER EDITIEREN: 
% Typ der Arbeit (für Deckblatt und ehrenwörtliche Erklärung)
% - bitte Zutreffendes auswählen
%\newcommand{\typMeinerArbeit}{1. Projektarbeit} 
%\newcommand{\typMeinerArbeit}{2. Projektarbeit} 
%\newcommand{\typMeinerArbeit}{Seminararbeit} 
\newcommand{\typMeinerArbeit}{Bachelor thesis} 

% Thema der Arbeit (für ehrenwörtliche Erklärung)
% HIER EDITIEREN: 
\newcommand{\themaMeinerArbeit}{A design science approach to visualize blockchain transactions (for teaching purposes)}

\thispagestyle{empty}

\begin{spacing}{1}
\begin{center}	
~\vspace{0mm}

% HIER EDITIEREN: Titel der Arbeit
{\sffamily
\LARGE  
% \Large  % bei sehr langen Titeln ggf. etwas kleinere Schriftart wählen
\textbf{A design science approach to visualize blockchain transactions (for teaching purposes)}

%\bigskip
%\textbf{ggf. etwas länger}
}


\vspace{15mm}

% HIER EDITIEREN
{\Large \typMeinerArbeit}

\vspace{1cm}

% HIER ggf. EDITIEREN
submitted on May 07, 2018
\vspace{15mm}
Faculty of economics
\medskip

Business informatics degree programme
\medskip

% HIER EDITIEREN: Kurs eintragen
Course WWI2015F

\vspace{10mm}

by

\vspace{10mm}

% HIER EDITIEREN: Vorname und Name eintragen 
{\large\textsc{Lisa Martiny}}

\vspace{10mm}
\end{center}

\vfill

% HIER EDITIEREN
\begin{tabular}{ll}
Responsible person in the training center: & DHBW Stuttgart: \\
\hspace{0.4\linewidth} & \\
HPE & Professor, John Sören Pettersson \\
Ralph Beckmann  &  \\
Blockchain Advisor TS Consulting \\
\\
Sinature \\
\end{tabular}


\vspace{1cm}
%(etwas Platz für die Unterschrift des Betreuers aus der Ausbildungsstätte)
\end{spacing}

% falls ein Sperrvermerk erforderlich ist: 
%\begin{center}
%\small
%\textbf{Sperrvermerk}:
%Der Inhalt dieser Arbeit darf weder als Ganzes noch in Auszügen \\
%Personen außerhalb des Prüfungsprozesses und des Evaluationsverfahrens zugänglich gemacht werden, sofern keine anders lautende Genehmigung der Ausbildungsstätte vorliegt. 
%\end{center}


% HIER EDITIEREN: Meta-Daten für PDF-Datei
\hypersetup{pdftitle={Titel der Arbeit}}
\hypersetup{pdfauthor={Ihr Name}}
\hypersetup{pdfsubject={Bachelorarbeit DHBW Stuttgart 2015}}

%%% Umstellung der Seiten-Nummerierung auf i, ii, iii ... %%%
\pagenumbering{Roman} 

%%% Abstract (optionale Kurzfassung Ihrer Arbeit) %%%
\begin{abstract}
\thispagestyle{kapitelkopfzeile}
\textbf{\LaTeX-Vorlage für Projekt-, Seminar- und Bachelorarbeiten}

Bei dem vorliegenden Dokument handelt es sich um eine Vorlage, die
für Projekt-, Seminar- und Bachelorarbeiten im Studiengang
Wirtschaftsinformatik der DHBW Stuttgart verwendet werden kann.

Sie setzt die technischen Vorgaben der Zitierrichtlinie des Studiengangs
(Stand: 01/2018) um.

\emph{Hinweis:} Dieses Dokument ersetzt keine Anleitung oder Einführung in \LaTeX,
für die Nutzung sind daher gewisse Vorkenntnisse unerlässlich. Ein Einstieg in 
\LaTeX\ ist aber weniger schwierig, als es vielleicht auf den ersten Blick scheint
und lohnt sich für das Verfassen wissenschaftlicher Arbeiten in jedem Fall.\footnote{%
so auch \url{http://www.spiegel.de/netzwelt/tech/textsatz-keine-angst-vor-latex-a-549509.html}} 

Ihre Rückmeldungen und Anregungen zu dieser Vorlage nehme ich gerne per E-Mail an die Adresse
\url{tobias.straub@dhbw-stuttgart.de} entgegen.

--- Prof. Dr. Tobias Straub

\vspace{5em}

\begin{center}\small
\begin{tabular}{ccl}
\multicolumn{3}{c}{\textbf{Versionshistorie}}\\
\hline
1.0	& 05.02.2015 & erste Fassung \\
\hline
1.1 & 16.02.2015 & siehe~\ref{anhang:ReleaseNotes11} \\
\hline
1.2 & 20.04.2015 & siehe~\ref{anhang:ReleaseNotes12} \\
\hline
1.3 & 20.02.2016 & siehe~\ref{anhang:ReleaseNotes13} \\
\hline
1.4 & 24.07.2017 & siehe~\ref{anhang:ReleaseNotes14} \\
\hline
1.5 & 07.01.2018 & siehe~\ref{anhang:ReleaseNotes15} \\
\end{tabular}
\end{center}

\end{abstract}


\cleardoublepage

\renewcommand{\listfigurename}{List of figures}
\renewcommand{\listtablename}{List of tables}

%%% Inhalts-, Abbildungs-, Tabellenverzeichnisse %%%
% sollen einzeilig gesetzt werden, um Platz zu sparen 
\begin{spacing}{1}
\tableofcontents
\clearpage
\chapter*{List of abbreviations}
\addcontentsline{toc}{chapter}{List of abbreviations}

Ein Abkürzungsverzeichnis ist optional. Das Paket \verb|acronym| kann weit mehr, als hier gezeigt.\footnote{siehe \url{http://ctan.org/pkg/acronym}}
Beachten Sie allerdings, dass Sie die Einträge selbst in sortierter Reihenfolge angeben müssen.

\begin{acronym}[DHBW] 
% Argument definiert die Breite der ersten Spalte anhand des längsten vorkommenden Eintrags
\acro{CRM}{Customer Relationship Management}
\acro{DHBW}{Duale Hochschule Baden-Württemberg}
\acro{IEEE}{Institute of Electrical and Electronics Engineers}
\acro{ITIL}{IT Infrastructure Library}
\acro{RoI}{Return-On-Invest}
\acro{UCS}{Universal Character Set}
\acro{UTF-8}{8-Bit UCS Transformation Format} 
\end{acronym}

\vspace{2em}
{\footnotesize
\textbf{Ergänzende Bemerkung:}
Eine im Text verwendete Abkürzung sollte bei ihrer ersten Verwendung erklärt werden. Falls Sie sich nicht selbst darum kümmern möchten, kann das das Paket \verb|acronym| übernehmen und auch automatisch Links zum Abkürzungsverzeichnis hinzufügen. Dazu ist an allen Stellen, an denen die Abkürzung vorkommt, \verb|\ac{ITIL}| zu schreiben. 

Das Ergebnis sieht wie folgt aus: 
\begin{itemize}
\item erstmalige Verwendung von \verb|\ac{ITIL}| ergibt: \ac{ITIL},
\item weitere Verwendung von \verb|\ac{ITIL}| ergibt: \ac{ITIL}
\end{itemize}
Wo benötigt, kann man mit dem Befehl \verb|\acl{ITIL}| wieder die Langfassung ausgeben lassen: \acl{ITIL}.

Falls man die Abkürzungen durchgängig so handhabt, kann man durch Paket-Optionen (in \verb|_dhbw_praeambel.tex|)
erreichen, dass im Abkürzungsverzeichnis nur die tatsächlich verwendeten Quellen aufgeführt werden (Option: \verb|printonlyused|) und zu jedem Eintrag die Seite der ersten Verwendung angegeben wird (Option: \verb|withpage|).
}
\clearpage
\thispagestyle{kapitelkopfzeile}
\listoffigures
\addcontentsline{toc}{chapter}{List of figures} % Abb.verz. ins Inh.verz. aufnehmen
\clearpage
\addcontentsline{toc}{chapter}{List of tables}   % Tab.verz. ins Inh.verz. aufnehmen
\listoftables
\end{spacing}

%%% Umstellung der Seiten-Nummerierung auf 1, 2, 3 ... %%%
\cleardoublepage
\pagenumbering{arabic}

%%% Ihr eigentlicher Inhalt %%%
% Empfehlung: strukturieren Sie Ihren Text in einzelnen Dateien 
% und binden Sie diese hier mit \input{includes/dateiname.tex} ein

\chapter{Einleitung}

Bald kann nun der Text Ihrer Projekt- oder Bachelorarbeit beginnen. Dank LaTeX wird Ihre Arbeit garantiert professionell aussehen. Für den Inhalt sind Sie aber weiterhin selbst verantwortlich~\mbox{;-)}

Natürlich ist es schwer, sich vorzustellen, wie das Dokument aussieht, wenn die Vorlage doch gar keinen Text enthält. Aus diesem Grund wird mit Hilfe des Pakets \enquote{blindtext} so genannter Blindtext erzeugt. Mit dem Befehl \verb|\blinddocument| wird nachfolgend ein ganzes Kapitel sinnfreier Blindtext eingefügt.\footnote{Beachten Sie, dass Sie in Ihrer Arbeit eine Strukturierung wie in Abschnitt 2.1 vermeiden sollten: Dort gibt es einen Abschnitt 2.1.1, aber keinen Abschnitt 2.1.2.} 

In den Abschnitten~\ref{section:werkzeuge} bis \ref{section:zeichencodierung} sind zuvor einige wichtige Einstellungen beschrieben. In Kapitel~\ref{chapter:zitate} finden sich Beispiele, wie Sie Quellen korrekt zitieren können. In Kapitel~\ref{chapter:abbildungenTabellen} werden Abbildungen, Tabellen, ein Code-Listing und auch mathematische Formeln in den Text eingebunden. Ab Seite~\pageref{chapter:quellen} finden Sie die Quellenverzeichnisse -- getrennt nach Artikeln/Büchern, Online-Quellen und Gesprächen. 

\emph{Hinweis:} Die farbigen, anklickbaren Links, die in der PDF-Ansicht enthalten sind, werden beim Ausdruck nicht wiedergegeben.\footnote{Das Feature lässt sich abschalten, indem man die Option \texttt{hidelinks} bei \texttt{documentclass} zu Beginn des Hauptdokuments hinzufügt.}

\section{Werkzeuge}\label{section:werkzeuge}

Sämtliche benötigten Werkzeuge sind Open Source und damit kostenlos nutzbar.

Für einen einfachen Start in \LaTeX\ sowie Tests können Sie einen Online-Editor wie \url{overleaf.com} verwenden. Allerdings sollten Sie eine Projekt- oder Bachelorarbeit nicht mit diesem System erstellen, da es sich um eine Cloud-Lösung handelt, bei der letztlich nicht gewährleistet ist, dass Firmen-Interna und schützenswerte Daten nicht in die Hände Dritter gelangen können. Für eine lokale Installation können Sie z.B.\ TeXLive (alle gängigen Plattformen, \url{http://tug.org/texlive/}), MikTeX (Windows, \url{http://www.miktex.org/}) oder TexShop (Mac OS, \url{http://pages.uoregon.edu/koch/texshop/}) verwenden.

Erstellen Sie Ihr Dokument von Beginn an in \LaTeX. Es ist etwa wenig sinnvoll, zuerst in Word zu schreiben und das Ergebnis am Ende nach \LaTeX\ zu konvertieren.\footnote{Falls es je nötig sein sollte, gibt es für die umgekehrte Richtung diverse Konverter, z.B.\ latex2rtf, \url{http://sourceforge.net/projects/latex2rtf/}.} 

   
Die Vorlage verwendet biblatex,\footnote{siehe etwa \url{http://www.ub.uni-konstanz.de/serviceangebote/literaturverwaltung/bibtex/bibtex-und-biblatex-benutzen.html} für eine Gegenüberstellung von BibTeX und biblatex.}
um den vorgegebenen Stil der Zitierrichtlinien umzusetzen. Der Editor muss daher so konfiguriert werden, dass er beim Übersetzen biber statt bibtex verwendet. Zur Verwaltung der Quellen nutzt man etwa JabRef (Java, \url{http://jabref.sourceforge.net/}).

\section{Benötigte Dateien}

Die folgenden Dateien, deren Namen einheitlich mit \verb|_dhbw_| beginnt, werden eingebunden. Änderungen an diesen Dateien sind nicht erforderlich.

\verb|_dhbw_authoryear.bbx       (Anpassung Einträge im Literaturverzeichnis)| \\
muss sich im selben Verzeichnis wie die zu kompilierende Hauptdatei (\verb|latex-vorlage.tex|) befinden, die folgenden Dateien im Unterverzeichnis \verb|includes|:
\begin{verbatim}
_dhbw_biblatex-config.tex  (weitere Einstellung für Biblatex)
_dhbw_erklaerung.tex       (ehrenwörtliche Erklärung)
_dhbw_kopfzeilen.tex       (Kapitelname in Kopfzeilen) 
_dhbw_praeambel.tex        (Einbindung der benötigten Pakete)
\end{verbatim}

Weiterhin sind nur für die Übersetzung der Beispieldatei erforderlich:
\begin{verbatim}
dhbw.png                   (Beispiel für eine Grafik, die eingebunden wird)
HelloWorld.java            (eingebundenes Java-Listing)
literatur-datenbank.bib    (Literatur-Datenbank mit Beispiel-Einträgen)
\end{verbatim}
Die Grafik ist im Verzeichnis \verb|graphics| abgelegt, die anderen Dateien in \verb|includes|. 

Ferner gliedert sich dieser Text in folgende Dateien (ebenfalls im Verzeichnis \verb|includes|), die per \lstinline[language=TeX]|\input|-Befehl eingebunden werden:
\begin{verbatim}
abbildungen_und_tabellen.tex
abkuerzungen.tex
abstract.tex
anhang.tex
deckblatt.tex
einleitung.tex
text_mit_zitaten.tex
\end{verbatim}

Um das Editieren (und Debuggen) zu erleichtern, ist es ratsam, ein längeres Dokument in einzelne Dateien zu strukturieren (z.B.\ kapitelweise).

\section{Vorzunehmende Einstellungen}

\emph{Hinweis:} Mit der ab 1/2016 geltenden Fassung der Zitierrichtlinien ist die zuvor bestehende Möglichkeit, auch doppelseitig auszudrucken, gestrichen worden. Gleiches gilt für die nun nicht mehr bestehende Option, \enquote{ebenda} zu nutzen.

\subsection{Deckblatt}

Die folgenden Anpassungen sind in der Datei \verb|deckblatt.tex| vorzunehmen. Die entsprechenden Stellen sind im Source Code wie folgt gekennzeichnet:
\lstset{language=TeX} 
\begin{lstlisting}
% HIER EDITIEREN: 
\end{lstlisting}

\subsubsection{Titel und ehrenwörtliche Erklärung}
Natürlich ist das Deckblatt anzupassen, schließlich soll dort Ihr Name erscheinen. Denken Sie bitte aber auch unbedingt daran, am Anfang der Datei festzulegen, um welchen Typ (Projekt-/Bachelor-/Seminararbeit) es sich handelt und wie Ihr Thema (Titel der Arbeit) lautet.\footnote{%
Falls Sie einen langen Titel mit Untertitel haben, können Sie für die Aufnahme in die ehrenwörtliche Erklärung den Untertitel auch weglassen.} Beide Angaben werden automatisch in die ehrenwörtliche Erklärung eingefügt, so dass Sie in der Datei \verb|_dhbw_erklaerung.tex| keine Änderungen vornehmen müssen.
In der Datei \verb|deckblatt.tex| wird die Schriftgröße für den Titel auf \verb|\LARGE| eingestellt. Falls Sie je einen sehr langen Titel haben, wählen Sie eine kleinere Schriftgröße mit dem Befehl \verb|\Large| (Zeile 21 im Source Code).

\subsubsection{Meta-Daten im PDF}
Am Ende von \verb|deckblatt.tex| können Sie Ihren Namen und den Titel der Arbeit für die Meta-Daten der PDF-Datei angeben. 


\subsubsection{Sperrvermerk}
Ein so genannter \enquote{Sperrvermerk} sollte eher die Ausnahme sein. Meiner Erfahrung nach enthalten die allerwenigsten Arbeiten brisante Firmengeheimnisse. Arbeiten werden -- unabhängig von einem Sperrvermerk -- von der DHBW sowieso nicht an Dritte weitergegeben, auch die Prüfer müssen die Inhalte vertraulich behandeln. Durch einen Sperrvermerk schränken Sie sich aber möglicherweise selbst ein, da Sie dann Ihre Arbeit im Grunde auch niemandem ohne Genehmigung Ihrer Firma zeigen dürfen.
Falls ein Sperrvermerk erforderlich ist, können Sie die auskommentierten Zeilen am Ende der Datei \verb|deckblatt.tex| nutzen.

\subsection{Hauptdokument}

Sie bearbeiten folgenden Abschnitt, um Ihre einzelnen Kapitel einzubinden:
\lstset{language=}
\lstinputlisting[firstline=42,lastline=53,firstnumber=42]{latex-vorlage.tex}

Je nachdem, ob Sie Abstract und Abkürzungsverzeichnis verwenden, ändern Sie die folgenden Zeilen.
In aller Regel werden Sie auch Tabellen und Abbildungen nutzen, ansonsten entfernen Sie die entsprechenden Verzeichnisse. 
\lstinputlisting[firstline=19,lastline=36,firstnumber=19]{latex-vorlage.tex}

Änderungen außerhalb dieser Bereiche sind nicht erforderlich.

\section{Zeichencodierung}\label{section:zeichencodierung}

Achten Sie darauf, dass Sie für alle \LaTeX\ und BibTeX-Dateien eine einheitliche Zeichencodierung verwenden, damit Umlaute und Sonderzeichen korrekt wiedergegeben werden. Sie müssen dazu ggf.\ die Einstellungen Ihres Editors anpassen.

Damit Sonderzeichen korrekt dargestellt werden, sollte als Codierung \emph{UTF-8} (Unicode) eingestellt sein. UTF-8 stimmt in den ersten 128 Unicode-Zeichen mit dem ASCII-Zeichensatz überein, kann aber auch Sonderzeichen oder Zeichen beliebiger Sprachen darstellen.

Diese Datei ist \emph{UTF-8}-codiert, weshalb sich in der Präambel folgender Befehl findet:
\lstset{language=TeX} 
\begin{lstlisting}
\usepackage[utf8]{inputenc}   % Zeichencodierung UTF-8 für Eingabe-Dateien
\end{lstlisting}

Mit der richtigen Codierung können die Sonderzeichen ä, ö, ü, Ä, Ö, Ü, ß, € wie gewohnt direkt im Source Code eines \LaTeX-Dokuments geschrieben werden.

\section{Fehlerbehebung}

\subsection{Kopfzeile}
Es ist \emph{kein Fehler}, sondern beabsichtigt, dass auf der ersten Seite eines Kapitels (und auch bei der ersten Seite der Verzeichnisse) die Kopfzeile fehlt.

\subsection{Kontrolle der Seitenränder}
Wenn Sie in Zeile 3 der Hauptdatei das Packet \enquote{showframe} einbinden, können Sie sich die Ränder Ihres Dokuments zur Kontrolle anzeigen lassen. Es empfiehlt sich, die Ränder zu kontrollieren und darauf zu achten, dass beim Druck \emph{keine Skalierung} im PDF-Viewer gewählt wird.  

Überlange Zeilen erkennen Sie außerdem an den Warnmeldungen \verb|Overfull \hbox| in der \verb|.log|-Datei.

\subsection{biber}
Zumindest unter Mac OS kommt es manchmal vor, dass sich biber aufhängt und mit einer Fehlermeldung folgender Form stoppt:

{\small
\begin{verbatim}
read_file '/var/folders/ay/ay9RQK7FEcKOxehY75+N4k+++TI/-Tmp-/
par-746f62696173737472617562/cache-a3cdad92316c60c9c5179d80d6bb51a7a024393c/
inc/lib/Biber/LaTeX/recode_data.xml' - sysopen: No such file or directory at 
/var/folders/ay/ay9RQK7FEcKOxehY75+N4k+++TI/-Tmp-/par-746f62696173737472617562/
cache-a3cdad92316c60c9c5179d80d6bb51a7a024393c/inc/lib/Biber/LaTeX/Recode.pm 
line 112.
INFO - This is Biber 1.9
INFO - Logfile is 'latex-vorlage.blg'
\end{verbatim}
}

Abhilfe schafft Löschen des temporären Verzeichnisses, hier: \verb|par-746f62696173737472617562|, inklusive der Unterverzeichnisse.

Ein korrekter biber-Lauf sieht in etwa so aus:
{\small
\begin{verbatim}
INFO - Logfile is 'latex-vorlage.blg'
INFO - Reading 'latex-vorlage.bcf'
INFO - Found 25 citekeys in bib section 0
INFO - Processing section 0
INFO - Looking for bibtex format file 'includes/literatur-datenbank.bib' for section 0
INFO - Decoding LaTeX character macros into UTF-8
INFO - Found BibTeX data source 'includes/literatur-datenbank.bib'
INFO - Overriding locale 'de-DE' defaults 'variable = shifted' with 'variable = 
non-ignorable'
INFO - Overriding locale 'de-DE' defaults 'normalization = NFD' with 
'normalization = prenormalized'
INFO - Sorting list 'nyt' of type 'entry' with scheme 'nyt' and locale 'de-DE'
INFO - No sort tailoring available for locale 'de-DE'
INFO - Writing 'latex-vorlage.bbl' with encoding 'UTF-8'
INFO - Output to latex-vorlage.bbl 
\end{verbatim}
}

\subsection{Zeilenumbruch bei langen URLs}
Sehr lange und komplizierte URLs\footcite{langeURL} können in manchen Fällen ein Problem beim Zeilenumbruch in Blocksatz darstellen, v.a.\ im Verzeichnis der Internetquellen. Folgende URL ragt beispielsweise in den rechten Rand hinein: 

\url{http://www.google.de/search?hl=de&source=hp&q=biblatex+umbruch+url&gbv=2&oq=biblatex+umbruch+&gs_l=heirloom-hp.3.0.0i13i30l2j0i22i10i30.1757.7464.0.8525.22.19.0.3.3.0.204.1959.13j5j1.19.0.msedr...0...1ac.1.34.heirloom-hp..0.22.1990.yWySyKFfLPY}

Eine nahe liegende Lösungsmöglichkeit ist, den betroffenen Satz etwas umzuformulieren.
Jetzt geht es allerdings auch nicht besser.\footnote{%
\LaTeX\ sieht von einer Trennung der URL bei einem Bindestrich ab, um Missverständnisse zu vermeiden.
}
\url{http://www.google.de/search?hl=de&source=hp&q=biblatex+umbruch+url&gbv=2&oq=biblatex+umbruch+&gs_l=heirloom-hp.3.0.0i13i30l2j0i22i10i30.1757.7464.0.8525.22.19.0.3.3.0.204.1959.13j5j1.19.0.msedr...0...1ac.1.34.heirloom-hp..0.22.1990.yWySyKFfLPY}

Man kann sich aber damit behelfen, dass man die URL manuell in mehrere hintereinander gestellte \verb|\href|-Befehle auftrennt (der anklickbare Link bleibt so weiterhin korrekt). Einzig die Zeichen \verb|&| und \verb|_| sind zu escapen (als \verb|\&| bzw.\ \verb|\_|). Schauen Sie bitte in den Source Code.

\newcommand{\mylongurl}[1]{%
% erstes Argument von href ist das Linkziel, das beim Klicken angesprungen wird
\href{http://www.google.de/search?hl=de&source=hp&q=biblatex+umbruch+url&gbv=2&oq=biblatex+umbruch+&gs_l=heirloom-hp.3.0.0i13i30l2j0i22i10i30.1757.7464.0.8525.22.19.0.3.3.0.204.1959.13j5j1.19.0.msedr...0...1ac.1.34.heirloom-hp..0.22.1990.yWySyKFfLPY}{\texttt{#1}}\\}

%beachten Sie, dass & und _ escaped werden müssen
\mylongurl{http://www.google.de/search?hl=de\&source=hp\&q=biblatex+umbruch+url\&gbv=2\&oq=}
\mylongurl{biblatex+umbruch+\&gs\_l=heirloom-hp.3.0.0i13i30l2j0i22i10i30.1757.7464.0.8525.22.}
\mylongurl{19.0.3.3.0.204.1959.13j5j1.19.0.msedr...0...1ac.1.34.heirloom-hp..0.22.1990.}
\mylongurl{yWySyKFfLPY}

Falls das Problem im Verzeichnis der Internetquellen auftritt, kann man sein Glück im Hauptdokument probieren mit:
\begin{lstlisting}
\sloppy
\printbibliography[heading=www,type=online]
\fussy
\end{lstlisting}


\textbf{Warnung:} Von der Verwendung von \verb|\sloppy| wird explizit abgeraten,\footnote{\url{http://mirrors.ctan.org/info/l2tabu/german/l2tabu.pdf}, Kategorie: \enquote{Todsünden}}
Seien Sie vorsichtig und prüfen Sie die Auswirkungen auf Ihr Dokument, falls Sie es doch tun!
\blinddocument
\chapter{Zitieren}\label{chapter:zitate}

Der Zitierstil ist so angepasst, dass er den Zitierrichtlinien des Studiengangs Wirtschaftsinformatik der DHBW Stuttgart entspricht. 

\section{Zitate in den Text einfügen}
In \LaTeX\ wird mit den Befehlen \verb|\footcite| 
oder 
\verb|\cite|
eine Referenz im Text eingefügt. Meist wird \verb|\cite| nur \emph{innerhalb} einer Fußnote benutzt. 
Damit ein vorangestelltes \enquote{vgl.} in der Fußnote erscheint, können Sie wie folgt zitieren:
\begin{verbatim}
\footcite[vgl.][S. 3]{Autor}
\footcite[vgl.][]{Autor}
\end{verbatim}

Das erste optionale Argument von \verb|\footcite| wird dem Zitat vorangestellt, das zweite ist die Seitenzahl. Den selben Effekt hätte
\begin{verbatim}
\footnote{vgl. \cite[S. 3]{Autor}}
\footnote{vgl. \cite{Autor}}
\end{verbatim}

Hinweis: Falls \enquote{vgl.}, aber keine Seitenzahl angeben werden soll, muss das zweite Argument vorhanden (jedoch leer) sein, ansonsten wird \enquote{vgl.} als Seitenzahl interpretiert. Falsch ist also: 
\begin{verbatim}
\footcite[vgl.]{Autor}  % so nicht!
\end{verbatim}



\section{Beispiele}
Nachfolgend ein paar Beispiele, um die korrekte Darstellung zu überprüfen:

\begin{itemize}
\item \cite{Schlosser} ist ein Buch über LaTeX.
\item Zur Vorlesung \emph{Logik und Algebra} gibt es das gleichnamige Lehrbuch.\footcite{Staab}
\item nochmal dasselbe Buch\footcite{Staab}
\item ein weiteres Buch desselben Autors\footcite{BuschlingerStaab}

\item Der Konferenzbeitrag \cite{Ancuti} beschäftigt sich mit Bildverarbeitung.
\item Cloud Computing wird in einer Diplomarbeit erklärt.\footcite[S.~14]{Boettger:Diplomarbeit}
\item Preiß\footcite{Preiss} gibt eine Einführung in Datenbanken.
\item Eine Erläuterung, was \enquote{Intangibles} sind, findet sich bei Stoi\footcite[S.~82]{Stoi}.
\item weitere Ausführung in derselben Quelle\footcite[vgl.][S.~84]{Stoi}
\item Laut Wikipedia\footcite[Abschnitt~5]{wiki:Wirtschaftsinformatik} ist Wirtschaftsinformatik ein interessantes Studienfach.
\item ITIL-Prozesse kann man tatsächlich auch mit \LaTeX\ dokumentieren.\footcite{Carvalho:PJ:2012-1}
\item Open-Source und Cloud-Computing in einem Buchbeitrag\footcite{Wind}
\item Buch mit zwei Autoren\footcite{MitZweiAutoren}
\item Buch mit drei Autoren\footcite{MitDreiAutoren}
\item Buch ohne Autor\footcite{OhneAutoren}
\item Buch ohne Autor und ohne Jahr\footcite{OhneAutorenOhneJahr}
\item und noch ein anderes Buch ohne Autor und ohne Jahr\footcite{OhneAutorenOhneJahr2}
\item Buch ohne Autor, aber dafür mit Herausgeber\footcite{keinAutorAberHerausgeber}

\item manche Bachelorarbeit baut auf einer vorhergehenden Projektarbeit\footcite{mayer:PA1} auf
\item das Handbuch zu BibLaTeX\footcite{biblatex:manual} und eines zu Windows 8\footcite{Win8}
\item zwei Beiträge zu Büchern\footcite{Trautwein:Nokia}${}^{,}$\footcite{Mann} und zu einem Konferenzband\footcite{Trautwein:Erfolgsfaktoren}
\item eine Online-Quelle\footcite{SAP:HANA}
\item ein Gespräch\footcite{Mueller:Gespraech}
\item eine plagiierte Dissertation,\footcite{GuttenPlag} nicht zur Nachahmung empfohlen
\item zum Testen, ob Umlaute und Sonderzeichen korrekt wiedergegeben werden\footcite{Umlauttest}
\item \textbf{Eindeutigkeit:} Normalerweise wird kein Vorname des Autors angegeben. Falls es allerdings zur Eindeutigkeit\footnote{\cite{trautwein2011unternehmensplanspiele} vs. \cite{hitzler2011optimierung}} (bei gleicher Jahreszahl) erforderlich ist, wird der Vorname abgekürzt bzw.\ nötigenfalls sogar ganz ausgeschrieben mit angegeben.\footnote{vgl. \cite{mayer:PA1} und \cite{mayerLukas:PA1}}
 
Welch ein Glück, dass Sie sich darum dank \LaTeX\ gar nicht kümmern müssen (arme Word\texttrademark-User ;-).

\item Verwendung von Sekundärliteratur\footnote{\cite{Primaerquelle} zitiert nach \cite{Sekundaerquelle}}
\end{itemize}


\section{Eintragstypen für die Literatur-Datenbank}

Die verwendete Literatur pflegen Sie in einer Literatur-Datenbank im Bibtex-Format. Dabei handelt es sich um eine Textdatei, wobei für jede Quelle mittels Name-Value-Pairs die relevanten Attribute (Autor, Titel etc.) hinterlegt sind. Die Datei wird üblicherweise nicht im Texteditor, sondern in einem spezialisierten Programm wie JabRef bearbeitet.

Sofern in der Literatur-Datenbank der Typ eines Eintrags (Entry Type) korrekt festgelegt ist, wird er im Literaturverzeichnis automatisch richtig dargestellt. Mit folgenden Typen sollten Sie i.d.R.\ auskommen:
\begin{description}
\item[article] Artikel in einer Fachzeitschrift, auch E-Journal (Zeitschrift in elektronischer Form)\footnote{Bei E-Journals/E-Books werden beim Zitieren anstelle der (u.U. nicht eindeutigen, da von der Schriftgröße abhängigen) Seitenzahl Abschnitt und Absatz näher bezeichnet, also: \cite[Abschnitt 1.2.3, Absatz 4]{Staab}.}
\item[book] Buch, auch E-Book 
\item[inbook] Kapitel in einem Buch, zu dem mehrere Autoren beigetragen haben 
\item[inproceedings] Beitrag zu einer Fachtagung/Konferenz 
\item[manual] Handbuch
\item[misc] anderweitig nicht zuordenbarer Typ, darunter fallen auch Gespräche
\item[phdthesis] Dissertation
\item[thesis] Bachelor-/Master-/Diplomarbeit (Art wird im Attribut \enquote{type} festgelegt) 
\item[online] Internet- oder Intranet-Quelle\footnote{\label{fn:onlineEntryType}Man beachte, dass der Eintragstyp \enquote{online} in JabRef einmalig eingerichtet werden muss.}
\item[report] technischer Bericht, Forschungsbericht oder White Paper; diesen Typ können Sie auch verwenden, um eine Projektarbeit zu zitieren (Art wird im Attribut \enquote{type} festgelegt) 
\end{description} 

Eine Übersicht über die notwendigen Attribute jedes Eintragstyps gibt die folgende Tabelle, wobei ein Schrägstrich als \enquote{oder} zu verstehen ist.\footnote{Auszugsweise entnommen aus \cite{biblatex:manual}.} Zudem sind die wichtigsten optionalen Attribute aufgeführt.

\begin{tabular}{l|p{6cm}|p{6cm}}
\textbf{Eintragstyp} & \textbf{notwendige Attribute} & \textbf{optionale Attribute (Auswahl)} \\
\hline
article & author, title, journal, year/date  & volume, number, pages, month, note \\
\hline
book & author, title, year/date  & publisher, edition, editor, \mbox{volume/number}, series, isbn, url \\
\hline
inbook & author, title, booktitle, year/date  & bookauthor, editor, volume/number, series, isbn, url \\
\hline
inproceedings & author, title, booktitle, year/date  & organization/publisher, editor, volume/number, series, isbn, url  \\
\hline
manual & author/editor, title, year/date  & organization/publisher, address, edition, month, note, url, urldate \\
\hline
misc & author/editor, title, year/date & howpublished, organization, month, note \\
\hline
phdthesis & author, title, institution, year/date  & address, month, note \\
\hline
thesis & author, title, institution, type, year/date &  address, month, note \\
\hline
online\textsuperscript{\ref{fn:onlineEntryType}} & author/editor, title, year\footnotemark/date, url & urldate \\
\hline
report & author, title, institution, type, year/date & number, version, url, urldate \\
\end{tabular}
\footnotetext{%
Sofern kein Jahr bekannt ist, sollte das Attribut nicht leer gelassen werden (sonst wird die aktuelle Jahreszahl automatisch eingefügt), sondern der Eintrag \enquote{o.J.} gewählt werden.}



\section{Mehrere Quellverzeichnisse}

\LaTeX\ ist in der Lage, mehrere Verzeichnisse für Quellen zu verwalten. Eine Zuordnung findet automatisch anhand des Eintragstyps in der Literatur-Datenbank gemäß folgender Regeln statt.\footnote{Falls Sie wissen möchten, wie ein separates Verzeichnis für Internet-/Intranet-Quellen erstellt werden kann, schauen Sie bitte in der Version 1.4 dieses Dokuments nach.}

\begin{itemize}
\item Ins Gesprächsverzeichnis werden alle Einträge vom Typ \enquote{misc} aufgenommen, sofern Sie zusätzlich im Attribut \enquote{keyword} den Wert \enquote{gespraech} stehen haben. 

\item In das allgemeine Literaturverzeichnis werden alle Einträge aufgenommen, die nicht in die vorige Kategorie gehören.
\end{itemize} 

\section{Zitieren von Sekundärliteratur}
Gelegentlich lässt es sich nicht vermeiden, aus der Sekundärliteratur zu zitieren. Dies leistet der folgende Befehl.
\begin{verbatim}
\footnote{\cite{Primaerquelle} zitiert nach \cite{Sekundaerquelle}}
\end{verbatim}
Ins Literaturverzeichnis soll nur die Sekundärquelle aufgenommen werden. Dies wird dadurch erreicht, dass in der Literatur-Datenbank bei der Primärquelle im Attribut \enquote{keyword} der Wert \enquote{ausblenden} eintragen wird.

\chapter{Beispiele für Abbildungen und Tabellen}\label{chapter:abbildungenTabellen}

Hier finden Sie Beispiele für Abbildungen, Tabellen, Formelsatz und  Source Code.

\section{Abbildungen}
In diesem Abschnitt gibt die Abbildungen~\ref{abb:Logo2cmHoch} und~\ref{abb:Logo2cmBreit}, die beide das Logo der DHBW zeigen.

\begin{figure}[htb]
\centering
\includegraphics[height=2cm]{graphics/dhbw.png}
\caption[DHBW-Logo 2cm hoch]{DHBW-Logo 2cm hoch.\footnotemark}
\label{abb:Logo2cmHoch}
\end{figure}
\footnotetext{Mit Änderungen entnommen aus: \cite{OhneAutorenOhneJahr}}

\lstset{language=TeX, % hervorzuhebende Keywords definieren
  morekeywords={footnotetext,footnotemark,footcite,caption}
}

\emph{Spezialfall:} Sofern \emph{innerhalb} der Bezeichnung einer Abbildung eine Fußnote angegeben oder eine Quelle referenziert werden soll, geschieht dies nicht per \lstinline|\footnote| oder \lstinline
|\footcite|. Vielmehr sind die Befehle \lstinline|\footnotemark| und \lstinline|\footnotetext| zu verwenden und außerdem das optionale Argument für \lstinline|\caption| anzugeben (vgl.\ Source Code).

\begin{figure}[htb]
\centering
\includegraphics[width=2cm]{graphics/dhbw.png}
\caption[DHBW-Logo 2cm breit.]{DHBW-Logo 2cm breit. (Quelle: DHBW\footnotemark)}
\label{abb:Logo2cmBreit}
\end{figure}
\footnotetext{\url{www.dhbw.de}}



\section{Tabellen}

In diesem Abschnitt gibt es zwei Beispiel-Tabellen, nämlich auf Seite~\pageref{tab:BeispielTabelleKlein} und auf Seite~\pageref{tab:BeispielTabelleGroesser}.

\begin{table}[htb]
\centering
\begin{tabular}{lcr}
links & Mitte & rechts \\
\hline
Muster & Muster & Muster \\
\end{tabular}
\caption{Kleine Beispiel-Tabelle.}
\label{tab:BeispielTabelleKlein}
\end{table}

\begin{table}[htb]
\centering
\begin{tabular}{|l|l|c|l|r||l}
    \textbf{Spalte 1} & \textbf{Spalte 2} & \textbf{Spalte 3} & \textbf{Spalte 4} & \textbf{Spalte 5} & \textbf{Spalte 6} \\
    \hline
    a        & b          & c                & d        & e        & f        \\
    Test     & Test, Test & Test, Test, Test & ~        & ~        & ~        \\
    1        & 2          & 3                & 4        & 5        & 6        \\
\end{tabular}
\caption{Größere Beispiel-Tabelle.}
\label{tab:BeispielTabelleGroesser}
\end{table}

\section{Etwas Mathematik}

Eine abgesetzte Formel:
\[
  \int_a^b x^2 \: \mathrm{d} x = \frac{1}{3} (b^3 - a^3)
\]

Es ist $a^2+b^2 = c^2$ eine Formel im Text.

\section{Source Code}

Source Code-Blöcke können auf folgende Arten eingefügt werden:

\lstset{language=Java}

Direkt im \LaTeX-Source Code:
\begin{lstlisting}
if(1 > 0) {
  System.out.println("OK"); 
} else {
  System.out.println("merkwuerdig");
}
\end{lstlisting}

oder eingefügt aus einer externen Datei.
\lstinputlisting{includes/HelloWorld.java}

\chapter*{Appendix}
\addcontentsline{toc}{chapter}{Appendix}
\section*{List of appendices}
\vspace{-8em}
\listofanhang
\clearpage
\spezialkopfzeile{Appendix} % damit in der Kopfzeile das Wort "Anhang" angezeigt wird

\anhang{Interviews}\label{Interviews}

\lstset{language=TeX, % hervorzuhebende Keywords definieren
  morekeywords={anhang, anhangteil}
}


% Um den Anforderungen der Zitierrichtlinien nachzukommen, wird das Paket \verb|tocloft| verwendet. Jeder Anhang wird mit dem (neu definierten) Befehl \lstinline|\anhang{Bezeichnung}| begonnen, der insbesondere dafür sorgt, dass ein Eintrag im Anhangsverzeichnis erzeugt wird. Manchmal ist es wünschenswert, auch einen Anhang noch weiter zu unterteilen. Hierfür wurde der Befehl \lstinline|\anhangteil{Bezeichnung}| definiert.

% In~\ref{anhang:abbildung} finden Sie eine bekannte Abbildung und etwas Source Code in~\ref{anhang:sourcecode}. 

\anhangteil{Interview with Daniel Kaltenbach}\label{anhang:InterviewDK}
This is the transcription of the complete interview between Daniel Kaltenbach (Interviewee, DK) and the author (Interviewer, LM). The interview was conducted on 22th of March, 2018 from 9:00 to 9:30 in person in the HPE office in Böblingen. 

\begin{xlist}
     \item[LM] Ok, jetzt haben wir das Interview gestartet. Noch eine Sache ganz vorweg, ist es okay, wenn ich deinen Namen in der Bachelorarbeit nenne oder möchtest du es lieber anonym machen?
     \item[DK] Lass uns das am Ende klären.
     \item[LM] Das war immer so am Anfang, also bei den Ausführungen in den Büchern ist immer noch: Ja, das muss man einmal festhalten. Aber erstmal etwas über dich! Was ist denn deine Rolle in HPE?
     \item[DK] Ja, also mein Name for the record. Ich heiße Daniel Kaltenbach und bin bei HPE  verantwortlich für das Thema IOT. Ich bin für Deutschland, Österreich und Schweiz sowie in Russland der designierte IOT Verantwortliche der eben die Kooperationsstrategie dieses Themas vorantreibt und eine cross-funktionale Inkubationsorganisation führt um das Thema IOT für HPE greifbar zu machen, Portfolio, Partnerschaften und strategische Initiativen zu treiben, zu verantworten und nachher zu Marktkapitalisierung zu führen.
     \item[LM] Ok, und wer zählt so zu deinen Kunden? Oder hast du direkt Kontakt zu Kunden?
     \item[DK] Ich habe ganz viel direkten Kontakt mit Kunden. Ich habe aber keine direkten eigenen Kunden, also ich bin ja kein Account Manager oder so, weil ich vertrete das Thema im Namen der HPE und damit ist mein Kunde die Geographie Deutschland, Österreich, Schweiz und Russland.
     \item[LM] Okay, worauf ich eher damit hinaus wollte, so was was macht ihr denn für eure Kunden? Also welche Angebote und Dienstleistungen von IOT her bietet ihr an?
     \item[DK] Enorm unterschiedlich. Ich nutze ganz gern die Analogie, die die PDC geprägt hat. Wenn sie IOT definieren, dann geht es im Prinzip um intelligente Produkte: Die Kaffeemaschine, die nach hause telefoniert und weiß, wann sie kaputt gehen wird. Wohlgleich das zweite Thema des  intelligenten Prozesses, wo es eben um horizontale Integration von Prozessketten geht, die schlussendlich in der Summe intelligenter werden, oder nicht. Also wenn ich weiß, was eben nach mir passiert, dann kann ich möglicherweise nacharbeiten was davor nicht so gut gelaufen ist. Ich kann die Qualität im Gesamtprozess betrachten, den Durchlauf optimieren. Wenn ich  weiß, dass der vor mir fertig ist kann die nachfolgende Maschine schon einmal anlaufen und sich aufwärmen, weil sie ja gleich einen Prozess bekommen wird et cetera. Das kann aber auch mein Customer Journey sein, der durch ein Retail läuft und eben, wenn jemand rein läuft dann ist es wichtig, dass ein Mitarbeiter hinten weiß, da ist jemand reingelaufen, damit ich ihn möglicherweise richtig ansprechen kann. Solche Dimensionen, also intelligentes Produkt oder intelligenter Prozess und darin beraten wir und liefern jetzt die Technologien, die helfen diese Daten von diesen Informationen, die da sind, Herr zu werden und darauf basierend intelligente Entscheidungen zu treffen.
     \item[LM] Ok dann würde ich jetzt auch schon direkt zum nächsten Teil gehen. Zum Thema: das Problem beim Verständnis von Blockchain. Ganz allererst: Was verstehst du genau unter Blockchain und was stellst du dir unter der Funktionsweise davon vor?
     \item[DK] Ich stelle mir erstmal vor, dass es etwas unheimlich Kompliziertes ist und ich glaube ich nehme damit teilweise schon die Frage vorweg was das Problem dadran ist. Ich glaube das ist Teil des Problems. Also es ist eine Technologie ja erstmal, die eben ähnlich einer Datenbank funktioniert aber eben dezentral die Daten hält und gegenseitig autorisiert. Das würde ich jetzt mal so wie als eine Blockchain beschreiben, wie eine Art verteiltes Dateisystem, verteiltes Datenbanksystem, das eben aber ein in sich, durch die Verteilung notwendiges Problem löst, eben die Autorisierung der Daten und der Richtigkeit. Wahrscheinlich medium richtig, aber das war jetzt zumindest mein Verständnis.
     \item[LM] Das reicht ja auch und wie hast du dir… also woher hast du jetzt dieses Verständnis?Also hast du das von vom Hörensagen, von irgendwelchen Artikeln oder wie hat sich das so zusammen geflochten?
     \item[DK] Mischung von allem. Ich glaube, dass warum es noch so medium richtig ist, liegt daran dass ich mich einfach nie selbst mit dem Thema intensiv auseinandergesetzt habe und einfach mal wirklich hinter das Thema gestiegen bin und deshalb hier die Arbeit auch ganz gut, das ist was, was die Industrie an der Stelle momentan nicht schafft: Eine sinnvolle Artikulation der Technologie und damit verbunden der Anwendungsszenarien. Also die Übersetzungen von einfachen menschen, wie mir hin zu „ Wir verstehen die Technologie und wir benutzen sie für dich“. Dazu ist mein Wissen nicht ausreichend und das schafft die Industrie auch nicht. Also mein Wissen basiert auf unterschiedlichsten Gesprächen, Konferenzen und ich war durchaus mal bei dem ein oder anderen Blockchain-artigen Workshop, wo durchaus nochmal so Grundlagen waren, aber immer wenn ich es gehört habe war es immer irgendwie nochmal ein bisschen neu und ein bisschen anders. Und einer hatte das so und der andere wieder so und ich bin halt nicht Techniker, ich mache Vertrieb am Ende des Tages.
     \item[LM] Ja das ist auch je nachdem wo man im Internet guckt, gibt es ganz viele verschiedene Ansätze das zu erklären. Ich versuche eben auch einmal die zu konsolidieren und zu schauen, dass man einen Ansatz hat der das ganzheitlich aber auch relativ einfach erklärt. Und wieso noch einmal meintest du denn, dass du dich eben nicht weiter damit auseinandergesetzt hast? Also du meintest ja so: Ja es ist ein kompliziertes Thema; hat dich das einfach davon abgeschreckt oder war es die fehlende Zeit? Was waren so die Gründe, dass du jetzt noch nicht da groß hinter stehst. 
     \item[DK] Zeit und Relevanz. Also wenn ich mich mit jeder Technologie irgendwie auseinandersetzen würde, die irgendwie IOT relevant ist, dann würde ich wahrscheinlich die nächsten zehn Jahre studieren und hätte es dann erst recht nicht gelernt, weil in den zehn Jahren wieder neue Sachen draußen sind. Also ich muss natürlich extrem selektiv wahrnehmen und das tue ich auch. Ich verstehe und das ist glaube ich mehr dieses HörenSagen dass es eine signifikante Implikation haben kann für bestimmte Anwendungsszenarien. Ich habe bisher für mich noch nicht den Bedarf verspürt mich intensiver damit auseinanderzusetzen, weil immer dann wenn so eine Diskussion aufgekommen wäre oder hätte kommen können, habe ich mir einfach einen Experten dazu geholt. Das bringt meine Rolle glücklicherweise mit, dass ich darin nicht Experte sein muss.
     \item[LM] Ja du meintest gerade, dass es schon Implikationen für IOT/deinen Bereich hat. Hast du da irgendwelche Beispiele im kopf, wo du sagen könntest: Ja Blockchain könnte da und da bedeutend sein? Also eher für deinen Bereich, jetzt nicht auf Finanzderivate bezogen.
     \item[DK] Ja, du wirst jetzt unfair. Nenee, ich glaube es ist ein gute Frage! Ich habe da jetzt natürlich irgendwie 20 Beispiele im Kopf, die irgendwie so allgemein diskutiert werden. Aber wenn ich jetzt so darüber nachdenke, kann ich jetzt nicht mit Bestimmtheit sagen, dass ich weiß, dass dieser UseCase oder dieser Anwendungsfall wirklich idealtypisch für eine Blockchain-Implementierung ist, weil ich auch weiß Blockchain ist eine Technologie, ne. Es gibt auch noch mögliche andere Technologien, die möglicherweise diesen Use Case genauso gut oder besser gar bedienen können. Von daher kann ich dir jetzt gerade keinen idealtypischen Fall nennen, weil mir wahrscheinlich auch die Technologie zu unbekannt ist.
     \item[LM] Darauf möchte ich auch eigentlich gar nicht hinaus, also es soll jetzt nicht DER UseCase sein, wo du denkst, auf jeden Fall muss da Blockchain hin. Sondern es geht auch gerade um dieses HörenSagen, was die meisten anderen Leute so als UseCases nennen. Ich würde gerne von dir wissen, was du denn so dazu zählst.
     \item[DK] Ja, ähm, alles was natürlich in den Bereich Monetarisierung geht, wo man eben den Mediator ausschalten möchte, also wie eine Bank oder so, wo man ein aufwändiges Transaktionssystem  benötigen würde. Wo eben bei Maschinen mit Maschinen Payments irgendwie ablaufen, würde ich jetzt mal so zu diesen Daten-Monetarisierungsbereich und Maschinen zu Maschinen Payments und natürlich das Thema Smart Contracting in der zweiten Dimension. Wo eben Maschinen miteinander verhandeln und so sich schließlich einigen auf Arbeitsausführungen und ja dann das Payment möglicherweise als zusätzliche notwendige Komponente mit einfließen lassen aber nicht notwendigerweise. Das wären jetzt so die zwei schwergewichtigsten. Aber ich bin wahrscheinlich zu weit weg um sagen zu können, das wäre perfekt oder so.
     \item[LM] Das ist ja alles gut und musst du auch gar nicht. Gut, zu der Frage wieso das so schwer zu verstehen ist haben wir ja schon festgehalten, dass es eben eine komplizierte Technologie ist. Meinst du da gibt es noch andere Punkte, dass sich eben die Leute/die Personen/potenzielle Kunden in deinem Gebiet sich noch nicht so sehr mit dem Thema auseinandergesetzt haben oder dieses Halbwissen haben, woran das liegen könnte?
     \item[DK] Mehrere Antworten. Weiterführend zu dem Thema "komplizierte Technologie" ist glaube ich der typischerweise  (und ich glaube das ist in allen Technologien so) beginnen über die Technologie zu sprechen, dann ist das immer ein Technologen-Kreis. Und ein Mensch/ein Typ Mensch, der natürlich hochtechnisch ist, wenn der dir versucht nachher, also ich meine dir und mir, ich will dir jetzt nichts unterstellen, aber zumindest mir, der einfache Mensch, dann versucht etwas kompliziertes zu erklären dann muss der immer fachlich total richtig sein. Wenn Physiker, der dir gegenüber sitzt und dir die Welt erklärt. Die Welt ist manchmal ganz einfach. Einen Fuss vor den anderen setzen, bewegst dich. Ja das ist für den nicht so einfach, das sind Energie, Bewegungsabläufe, da sind Erdanziehung, schlag mich tot. So ist es hier auch ein bisschen. Ich glaube, es ist noch nicht die Community oder die Technologiegruppe nenne ich sie mal, jetzt mal irgendwie verdreht, dass sie noch nicht in der Lage  sind eine einfache Sprache zu finden, als sie sich selbst einfach noch so Technologie-verliebt sind und den Transfer in die einfache Menschenwelt nicht schaffen. Nummer 1, also die Gruppe der Menschen, die damit beschäftigt sind. Die zweite Dimension ist wir leben momentan, glaube ich jetzt in so einer Technologiespirale in der einfach enorm viele Techniken gerade unterwegs sind, die irgendwie alle auf uns einpressen. Sagen wir mal 5G als Kommunikationsstandard, das ist künstliche Intelligenz, 3D Druck, es gibt irgendwie 50 solcher Dinger. Und die menschen sind in Summe noch in so einer Abwartephase. Welche von den Dingern werden jetzt revolutionsartig meine Welt verändern, mit welchen muss ich mich auseinandersetzen? Das sind so viele, dass wir eine Überforderung haben an Dingen, die die Menschen typischerweise nicht verstehen, weil sie einfach kompliziert sind. Weil das eben die Evolution von Technik ist. Ich glaube, das sind so die zwei Dimensionen. Also, in der Gruppe der Technologen ist man noch nicht bereit eine einfache Sprache zu wählen, weil das technische Thema einfach noch zu technisch ist. Oder man sich aber so sehr die Technik differenziert. Und gleichwohl bewegen wir uns in einem Umfeld, wo Technologien so rapide da sind und auch wieder weggehen, dass die Leute sich nicht genug damit auseinandersetzen zu können.
     \item[LM] Vielen Dank. Jetzt zu der nächsten Frage, die war: So wir haben jetzt schon gesagt, ja welche möglichen Konzepte bei IOT für Blockchain gibt es. Denkst du auch, dass das für HPE und die Kunden interessant sein könnte, das mal einzuführen für ein paar UseCases?
     \item[DK] Technologie ist ein Vehikel. Und Blockchain ist an der Stelle ein Vehikel, weil es eine Datenbank ist, und Technik ist ein Vehikel, die am Ende des Tages Mehrwerte schaffen müssen.  Und für Kunden sind Mehrwerte relevant, wir müssen Business Outcomes mit denen diskutieren, Gründe, Verbesserungen des Produktes, für Qualität, wie waren wir beim Ausbringen und mehr ausbringen, Kosten reduzieren oder ganz klassisch und wissensfragen und sobald wir es schaffen zu übersetzen, warum Blockchain da einen Beitrag leisten kann oder ein auf Blockchain basierendes Anwendungsszenario, das deswegen komplizierter ist aber Blockchain als ein technologisches Element dadrin hat. Dann sind die hoch zufrieden damit und momentan sind wir noch in so einer Phase, wo die Leute geil finden, dass Blockchain da drin ist. Das keinen Mehrwert hat, weil es halt geil ist Blockchain gemacht zu haben. Der Sache wegen geil. Aber ich glaube, das kühlt sich in der nächsten, das dauert gar nicht mehr lange, in den nächsten sechs Monaten wird sich das noch ein bisschen regen und dann wird der Markt da sein, verstanden haben: Okay, ist es ein Werkzeug in einem großen Werkzeugkasten und setzt das ein bitte. Also ihr macht das, ihr löst das IT Problem da. Das ist offensichtlich eins der Werkzeuge, sowie beim Handwerker, mir ist dann vollkommen egal ob der einen großen Schraubenzieher oder kleinen hat, solange die Schraube eingedreht wird. Schraube in der Wand. Im Grunde: Reduktion von Kosten und wenn Blockchain das Thema dafür ist, dann geil machst du das. Aber ich glaube, das ist nicht momentan, sechs Monate noch.
     \item[LM] Glaubst du denn, dass wenn man dem Kunden dann erklärt wie die Blockchain funktioniert, ob das einen Einfluss darauf hat? Also so wie du es anläuten ließest war es ja eher nicht der Fall. Also dass er sagt, Blockchain für mich egal wie das funktioniert, Hauptsache es funktioniert. Oder denkst du auch, dass er daran interessiert sein könnte zu wissen, wie es funktioniert? 
     \item[DK] Mehrere Stufen. Zuerst interessiert ihn, was ist der Business Outcome. Und natürlich will er irgendwann verstehen, ist dieses Angebot, dass ich von dir bekomme, liebes Technologieunternehmen,  hat das Substanz? Wie ist das umgesetzt? Hat das, möglicherweise (weil wir erklären natürlich ihm die Straße nach nirgendwo oder zu seinem Business Problem). Aber er hat ein Ökosystem, in das er es einbetten muss. in ein Business Continuity System, das heißt er muss verstehen was passiert, wenn diese Lösung oder diese Technik von dir ausfällt. Gibt es ein Backup und Disasterplan, was wollen wir nicht. Und wir werden in diesem Angebot, in dieser Erstdiskussion, werden wir natürlich alle Aspekte erläutern. Das heißt in dem Moment, wo er eben tiefer einsteigen möchte oder wir bewiesen haben, ja wir können ihr technisches Problem lösen mit Technologie. Würde er fragen: Welche Technologie verwendest du und was sind dafür die Disasterprozesse? Was denn, wie sicher ist das? Ist die Technik morgen weg? Also gibt es ein Substitut?Was auch immer,  da wird es sicherlich eine Diskussion geben und dann ist die Technik wichtig, wenn dann die Technologie Blockchain eine entsprechende Kredibilität erreicht, die ihm nützt, dann absolut. Es gibt Technologien von der wir heute sagen würden, wir machen den Prozess und der Kunde guckt dahinter und du würdest sagen: Und ich liefere Ihnen dann eine DVD. Dann würde er tot umfallen. Wahrscheinlich wenn ich eine Kassette noch liefern würde, noch zweimal tot umfallen und sagen: Ihr seid kein kredibles Unternehmen, weill sie keine Technologien nehmen, die nicht mehr State of the art ist, und eben kein zukünftiges ordentliches Substitut mehr hat oder keine Wiedersupply hat. Von daher hat es da sicherlich eine Relevanz, wenn sich die Technik durch setzt, aber die Technologie alleine macht keinen Mehrwert.
     \item[LM] Ja, da stimme ich auch zu. Mhmm und woran meinst du liegt es, dass Blockchain momentan so ein Hype ist? Beziehungsweise ich empfinde das so,  weil ich mich mehr damit auseinandersetze, aber AI ist ja auch ein großes Thema. Woran glaubst du liegt das?
     \item[DK] Schwierig, ich glaube da gibt es jetzt wieder 1000 Theorien oder wir können über die mediale Diskussion darüber reden,  so Fake News. Es kann sein, dass es jemand möchte, dass es ein total wichtiges Thema ist und gesteuert ist deswegen, Boah Verschwörungstheorien. Aber ich weiß es nicht, um ehrlich zu sein. Ich glaube durchaus, dass da eine Dimension drin ist, dass die Technik ziemlich gut ist.  Also ich glaube, das ist durchaus ein technologischer Durchbruch, der Eigenschaften mitbringt, die tatsächlich einige Anwendungsszenarien revolutionär sein können. Sonst hätte sich wahrscheinlich dieser Hype auch nicht so lange gehalten und zumindest bei uns in der Technologie-Community würde sich das nicht so lange halten, aber ich glaube da sind wir zu kritisch als Industrie. Von daher gehe ich mal schwer davon aus, dass da ein revolutionsartiges Element drin ist, das eben ganz, ganz besonders ist und uns deswegen weit nach vorne bringen kann. Ich kann jetzt auch nicht sagen, dass es genau diese eine Sache ist, die uns nach vorne gebracht hat. 
     \item[LM] Okay, aber Dankeschön. Und dann kommen wir jetzt auch schon zum dritten Teil: Was sollte denn in dieser Anwendung drin sein? Wie stellst du dir eine optimale Erklärung von Blockchain vor? Auf welchem Sprachniveau, weil wir das eben schon angeschnitten haben, und welche Konzepte?
     \item[DK] Also ich halte, im Grunde jetzt geht es um die Lerntheorie. Ich halte interaktive lernetheoretische Grundlagen als Basis. Ich glaube das sind die, die möglichst viele Sinne stimulieren, weil die für die Neulinge gut sind. Deshalb sind diese Erklärvideos immer sehr gut, wenn die dann auch eine Überleitung haben in einen praktischen Teil. So okay, jetzt mach das mal selbst. Oder für Erfahrene irgendwie selbst, hat das den höchsten Mehrwert. Ich glaube, weil es so ein hochtechnisches Thema ist, müssen wir unterschiedliche Ebenen bedienen und wir müssen das wahrscheinlich, halt mich daran jetzt nicht fest, in drei Ebenen diskutieren. Das sollte auf einer sehr abstrakten Art und Weise sein, dann muss es was eher technologisch vergleichend sein. Also mann muss sagen wie grenze ich mich gegen Technik links von mir und rechts von mir ab.  Keine Ahnung, welche die jetzt sind. Vielleicht Datenbank, etwas anderes,  verteiltes Dateisystem oder keine Ahnung, weiß ich jetzt nicht. Abgrenzen. das was nicht ist und dann muss es natürlich eine Erklärung der Technik zu der technischen Art und Weise haben, die einfach ausreichend einfach ist. Also nicht ein 10 Seiten Papier mit 0en und 1en, wo ich das erstmal mehr in Sprachtext übersetzen muss, sondern, dass - sage ich mal - in einer Art und Weise konsumierbar ist, dass ich wenn ich eine Viertelstunde aufgeregt lese und interessiert lese und dann sage: ja habe ich verstanden. Ich glaube, das sind mehrere Elemente, die wir schlussendlich brauchen, aber ich glaube dass, die größte Hürde bis heute ist der abstrakte Teil, das ganz oben, der möglicherweise auch diesen vergleichenden Teil drin hat. Weil das die Sprache ist, die mir noch fehlt. Ich habe mir eine Handvoll von diesen Erklärvideos angeschaut, immer mal wieder. Aber sie waren irgendwie nicht schlüssig, die waren nicht differenzierend genug,  und dass ich das Gefühl hatte, das erklärt mir jetzt gerade alles und nichts. 
     \item[LM] Gut, bei den nächsten Fragen wollte ich eigentlich darauf eingehen, welche Komponenten ich bei Blockchain. erklären sollte, aber von deinem Verständnis her... Blocks, weißt du, kennst du denn die Begriffe: Private Key, Public Key, Blocks, Hashfunktionen, kannst du dir darunter etwas vorstellen? Oder weißt du, dass das etwas mit Blockchain zu tun hat und hättest die Dinge auch gerne erklärt? Oder reicht da nur anschneidend?
     \item[DK] Ja genau das meine ich mit kaskadierend. Wir müssen, du musst das auf einer Ebene erklären, wo du diese Begrifflichkeiten nicht benötigst. Und natürlich, weil ich meine es gibt tonnenweise Informationen zu Public Key/Private Key Geschichten, aber ich glaube das musst du einweben. Und du musst eine Geschichte bauen, dass am Ende des Tages klar wird, was diese Technologie hat und das ist genau der Punkt: es nutzt diese Technologie (public und private key) für Security Geschichten, das ist jetzt nicht das einzige. Weil jetzt wieder ich als geneigter Nichtversteher sage mir: Gut, die haben jetzt das eine verwendet, aber wahrscheinlich gibt es da draußen 1000, und warum ist das jetzt genau das Richtige und warum ist das das Beste? Ist das substituierbar? Wenn jetzt so übermorgen jemand sagt, die Public Private Key Geschichte ist total hackbar, es ist falsch! Ist dann Blockchain per se auch falsch oder also es ist eine notwendige Verbindung oder eine Kann-Verbindung? Also das ist einfach zu viel für mich. Für mich, unpräzise oder unspezifisch, aber ich sage, weil das muss nur eben kaskadierend aufgebaut werden. Du musst erst ganz ganz oben anfangen. Erklären, was ist die Vision dahinter, was ist das Thema das du baust. Dann muss man eine Ebene tiefer gehen, also was sind die Bausteine technologischer Art und dann musst du erklären warum sind das genau diese und auch wieder abgrenzend zu links und rechts. Und dann sagen: Okay, jetzt erkläre ich es euch eine Ebene tiefer. Ich glaube dieses kaskadierende Modell baut Verständnis auf, in der jeder entscheiden kann, wann er austritt.  Für meine Rolle wäre es wahrscheinlich notwendig, alle Ebenen einmal durchgemacht zu haben. Sobald ich feststelle, dass es von relevanter Größe für die Industrie ist. Aber im Moment ist es, wen ich das Oberste haben möchte, dann würde ich die sprache von A, wenn ich das Unterste haben möchte die Sprache von Z und die zwei passen aber nicht zusammen. Also es gibt kein durchgängiges Erklärsystem, weil es einfach noch keiner gemacht hat, deshalb finde ich die Arbeit gut. 
     \item[LM] Dankeschön, ich hoffe das auch. Perfekt, die nächste Frage war nämlich: Wie sollte die Erklärung am besten aufgebaut sein? Das haben wir schon erledigt. Meinst du denn jetzt, wenn ich diese Erklärung habe von abstrakt runter und zurück... Noch einmal im Bereich Kundenprojekte. Ob man das dann dort auch anwenden sollte? Um zu sehen, also dass die Anwendung, die ich jetzt gestalte, auch beim Kunden gezeigt werden kann um mit ihm gemeinsam zu erarbeiten was Blockchain denn jetzt eigentlich bedeutet.
     \item[DK] Okay, ich habe mich jetzt gerade gefragt... Wenn ich jetzt wieder zwei Schritte zurücknehme, da müsste ich eigentlich sagen: Nein. Den Kunden interessiert Blockchain nicht, den Kunden interessieren Business Outcomes. So habe ich das gesagt. Aber du hast die Frage richtig gestellt. Um ein besseres Business Outcome mit ihm zu erarbeiten. Erarbeiten ist das wichtige dadrin.  Wir sind heute, glaube ich nicht in der Dimension oder nicht in der Marktsituation, wo wir ein "Friss"-Ergebnis hinlegen und der Kunde frisst am Ende des Tages. Wir sind nicht in dem Commodity-Markt und wir wissen, wenn ich ein Stückchen Hühnchen hole, dann frage ich nicht wo das herkam,  und wie es zerlegt wurde, weil ich gehe davon aus, dass es wenn es hinter dem Tresen liegt und den Bio-Stempel hat, dann war die Kühlkette sauber, was auch immer. So sind wir noch nicht. Das heißt, wir sind im komplexen Bereich. Also wir werden immer diese Ebene benötigen, wo ich erkläre und wir sind auch in einem sehr kooperativem Umfeld. Wir erarbeiten Lösungen mit Kunden und um eine Lösung zu erarbeiten müssen wir ihm eben Bausteine geben,  die er versteht. Und das ist genau meine Aufgabe. Und wenn wir da genau die richtige Sprache treffen. Wenn wir dem Kunden vermitteln, dass wir seine Sprache sprechen und das ist genau das, was wir hier bauen. Wir müssen seine Sprache sprechen um ihnen zu erklären, was der technologische Mehrwert dieser Technik ist um dabei ein Business Problem zu lösen. Wenn wir das sauber erklären können, dann bauen wir damit eine hohe Kredibilität auf, ein hohes Vertrauen und damit die Möglichkeit mit ihm was weitergehendes zu machen. Blockchain erklären kann theoretisch jeder, nur wir können es halt besser und wir treffen seine Sprache. Und damit sind wir der vertrauensvollere Anbieter wie irgendjemand anderes.
     \item[LM] Ok, perfekt, vielen Dank. Das war es auch schon. Wir sind tatsächlich sehr viel schneller durch gekommen, als ich gedacht habe, aber vielen Dank, du hast wirklich alles gut beantwortet und damit kann ich auch jetzt gut arbeiten.
     \item[DK] Abschließend: Ja, du darfst meinen Namen verwenden, sonst haben wir es wieder vergessen. 
\end{xlist}

\anhangteil{Interview with Ralph Beckmann}\label{anhang:InterviewRB}



\anhang{Release Notes}
\anhangteil{Änderungen in Version 1.1}\label{anhang:ReleaseNotes11}
In Version 1.1 sind einige Rückmeldungen, die nach der Einführungsvorlesung am 6.2.2015 oder nach Veröffentlichung der Vorlage in Moodle eingegangen sind, berücksichtigt worden. Korrekturen sind mit \enquote{(Fix)} gekennzeichnet. 

\begin{itemize}
\item \verb|latex-vorlage.tex|
\begin{itemize}
\item (Fix) Abkürzungsverzeichnis wird vor Abbildungsverzeichnis platziert
\item (Fix) Abbildungs- und Tabellenverzeichnis in Inhaltsverzeichnis aufgenommen
\item (Fix) Quellenverzeichnis wird nun ohne Kapitelnummer dargestellt

\item eingebundene Dateien in Unterverzeichnissen \verb|includes| bzw.\ \verb|graphics|
\item Beispiel-Anhang (Datei \verb|anhang.tex|) mit Erklärungen wurde eingebunden 
\end{itemize}

\item \verb|_dhbw_praeambel.tex|
\begin{itemize}
\item (Fix) das Paket hyperref wird nach biblatex eingebunden, um ein Problem mit der Verlinkung der Fußnoten im PDF zu beheben
\item (Fix) Fußnoten  gemäß der Richtlinien fortlaufend nummeriert und nicht pro Kapitel
\item Einstellungen hinzugefügt, um Anhangsverzeichnis zu ermöglichen
\item bessere Kompatibilität zwischen KOMA-Script (scrreprt) und anderen Paketen mittels scrhack
\end{itemize}

\item \verb|_dhbw_biblatex-config.tex|
\begin{itemize}
\item (Fix) keine Abschnittsnummern für einzelne Verzeichnisse im Quellenverzeichnis
\end{itemize}

\item \verb|abbildungen_und_tabellen.tex|
\begin{itemize}
\item Erklärung, wie eine Fußnote/ein Zitat bei einer Abbildung zu erstellen ist
\end{itemize}

\item \verb|abkuerzungen.tex|
\begin{itemize}
\item Abkürzungsverzeichnis wird im Inhaltsverzeichnis aufgeführt
\end{itemize}

\item \verb|abstract.tex|, \verb|anhang.tex|, \verb|einleitung.tex| 
\begin{itemize}
\item Erklärungen im Text ergänzt
\end{itemize}

\item \verb|deckblatt.tex|
\begin{itemize}
\item Meta-Daten (Autor, Titel) für die generierte PDF-Datei lassen sich nun festlegen
\end{itemize}

\end{itemize}


\anhangteil{Änderungen in Version 1.2}\label{anhang:ReleaseNotes12}
Über das Forum in Moodle sind einige Rückmeldungen eingegangen -- vielen Dank an alle, die dazu beigetragen haben. In der Version 1.2 wurden folgende Änderungen vorgenommen, wobei Korrekturen wieder mit \enquote{(Fix)} gekennzeichnet sind: 

\begin{itemize}
\item \verb|latex-vorlage.tex| (Hauptdokument)
\begin{itemize}
\item (Fix) Zeile 19: Seitenzahlen zu Beginn mit römischen \emph{Groß}buchstaben nummeriert
\end{itemize}

\item \verb|_dhbw_praeambel.tex|
\begin{itemize}
\item Zeile 39/40: Unterstützung für \enquote{ebenda} 
\item Zeile 46--68: zweite Gliederungsebene für Anhänge ermöglicht
\item (Fix) Zeile 70--73: Abbildungen und Tabellen: Zähler fortlaufend, kein Rücksetzen zu Kapitelbeginn (Paket \verb|chngcntr| anstelle von Paket \verb|remreset|)
\end{itemize}

\item \verb|_dhbw_biblatex-config.tex|
\begin{itemize}
\item (Fix) bei Quellen mit Herausgeber, aber ohne Autor wird der Name des Herausgebers im Verzeichnis fett gedruckt
\item Unterstützung für \enquote{ebenda} 
\end{itemize}

\item \verb|abkuerzungen.tex|
\begin{itemize}
\item Bemerkungen zur fortgeschrittenen Nutzung des \verb|acronym|-Pakets eingefügt 
\end{itemize}

\item \verb|einleitung.tex|
\begin{itemize}
\item Abschnitt 1.3 zu Einstellungen ergänzt
\item Abschnitt 1.5 zu Fehlerbehebungen eingefügt 
\end{itemize}

\item \verb|text-mit-zitaten.tex|
\begin{itemize}
\item Abschnitt 3.1 eingefügt, Erläuterungen zum Zitieren mit \enquote{vgl.} und \enquote{ebenda}. 
\item Abschnitt 3.2: Beispiele ergänzt
\item Hinweis zu Jahreszahlen bei Online-Quellen
\end{itemize}

\item \verb|anhang.tex|
\begin{itemize}
\item Erläuterungen zur zweiten Gliederungsebene
\end{itemize}

\item \verb|literatur-datenbank.bib|
\begin{itemize}
\item weitere Beispiele für Quellen
\end{itemize}

\end{itemize}

\anhangteil{Änderungen in Version 1.3}\label{anhang:ReleaseNotes13}
Durch die ab 1/2016 geltenden Änderungen der Zitierrichtlinien des Studiengangs waren einige kleinere Anpassungen der Vorlage erforderlich, die nachfolgend beschrieben sind. Bei dieser Gelegenheit ebenfalls erfolgte Korrekturen sind wieder mit \enquote{(Fix)} gekennzeichnet:

\begin{itemize}
\item \verb|latex-vorlage.tex| (Hauptdokument)
\begin{itemize}
\item Hinweis auf Option doppelseitiger Druck entfernt
\item Schriftgröße der Kapitelüberschriften verkleinert
\item (Fix) Kopf- und Fußzeilen werden nun korrekt angezeigt für erste Seite eines Kapitels und auch  Quellenverzeichnisse
\end{itemize}

\item \verb|_dhbw_praeambel.tex|
\begin{itemize}
\item Angabe des unteren Rands für Seitenzahl, da diese nun unten rechts steht
\item Unterstützung für \enquote{ebenda} entfernt
\item (Fix) Präfixe wie \enquote{von} im Namen eines Autors werden berücksichtigt
\item Anpassung der Abstände bei Kapitelüberschriften
\item Kopf- und Fußzeile für Verzeichnisse nun in \verb|_dhbw_kopfzeilen.tex| definiert 
\end{itemize}


\item \verb|deckblatt.tex|
\begin{itemize}
\item Schriftgröße des Titels vergrößert
\item Befehl \verb|\typMeinerArbeit| eingeführt, um Typ auszuwählen
\item Festlegung des Themas (für ehrenwörtliche Erklärung) mit Befehl \verb|\themaMeinerArbeit|
\item Darstellung der Angabe des Betreuers in der Ausbildungsstätte angepasst
\item Formulierung des Sperrvermerks angepasst  
\end{itemize}

\item \verb|_dhbw_erklaerung.tex|
\begin{itemize}
\item Formulierung angepasst an geänderte Prüfungsordnung
\item Typ und Thema der Arbeit werden automatisch eingefügt
\end{itemize}

\item \verb|_dhbw_kopfzeilen.tex|
\begin{itemize}
\item Seitennummern stehen jetzt unten rechts
\item (Fix) Kopf- und Fußzeile werden nun korrekt angezeigt in Verzeichnissen und dem Anhang
\end{itemize}

\item \verb|_dhbw_biblatex-config.tex|
\begin{itemize}
\item Anpassung des Zitierstils auf die ab 1/2016 geltenden Regelungen  
\item Vorkehrungen für Eindeutigkeit (Hinzufügen abgekürzter oder nötigenfalls ausgeschriebener Vorname) bei Übereinstimmung von Name und Jahreszahl 
\end{itemize}

\item \verb|einleitung.tex|
\begin{itemize}
\item Abschnitt 1.3 zu Einstellungen grundlegend überarbeitet
\item Abschnitt 1.5.2 zur Kontrolle der Seitenränder eingefügt 
\end{itemize}

\item \verb|text-mit-zitaten.tex|
\begin{itemize}
\item Abschnitt 3.1: Hinweise zu \enquote{ebenda} entfernt
\item Abschnitt 3.2: Beispiele zur Eindeutigkeit des Zitats ergänzt
\item Abschnitt 3.3: Hinweise für E-Journals/E-Books ergänzt 
\end{itemize}

\item \verb|anhang.tex|
\begin{itemize}
\item (Fix) Befehl \verb|\spezialkopfzeile| aufgenommen, damit in Kopfzeile das Wort \enquote{Anhang} angezeigt wird 
\item diese Release Notes wurden in eine eigene Datei verschoben
\end{itemize}

\item \verb|release_notes.tex|
\begin{itemize}
\item s.o.
\end{itemize}


\item \verb|literatur-datenbank.bib|
\begin{itemize}
\item weitere Beispiele für Quellen
\end{itemize}
\end{itemize}

\anhangteil{Änderungen in Version 1.4}\label{anhang:ReleaseNotes14}
Durch nicht abwärtskompatible Änderungen beim Versionswechsel von Biblatex 3.2 zu 3.3 sind einige Änderungen notwendig geworden.\footnote{Diese basieren auf Vorschlägen von Yannik Ehlert -- vielen Dank dafür!}
Die vorliegende Version 1.4 wurde erfolgreich mit MikTeX gestestet (portable Version 2.9.6361 vom 3.6.2017, unter Verwendung von Biblatex 3.7).

\begin{itemize}
\item \verb|_dhbw_biblatex-config.tex|
\begin{itemize}
\item Anpassung der \verb|\usebibmacro|-Befehle
\end{itemize}

\item \verb|_dhbw_authoryear.bbx|
\begin{itemize}
\item  Änderung von \verb|\printdateextralabel| zu \verb|\printlabeldateextra|
\end{itemize}
\end{itemize}

\anhangteil{Änderungen in Version 1.5}\label{anhang:ReleaseNotes15}
Für den Test dieser Version auf einem Windows-System wurde wieder die portable Version von MiKTeX (2.9.6521 vom 10.11.2017) verwendet.\footnote{\url{http://miktex.org/portable}} Da in diesem Paket leider die Versionen von Biblatex (3.10) und Biber (2.7) inkompatibel sind, ist es erforderlich, die Datei \verb|biber.exe| im Verzeichnis \verb|texmfs\install\miktex\bin\| durch die aktuelle Version 2.10 vom 20.12.2017\footnote{\url{https://sourceforge.net/projects/biblatex-biber/files/biblatex-biber/current/binaries/Windows/}} zu ersetzen. Im Editor TeXworks verwendet man dann zum Übersetzen des \LaTeX-Sourcecodes Typeset/pdfLaTeX bzw.\ Typeset/Biber.

Korrekturen sind wieder mit \enquote{(Fix)} gekennzeichnet.

\begin{itemize}
\item \verb|latex-vorlage.tex| (Hauptdokument)
\begin{itemize}
\item Nach der Änderung der Zitierrichtlinien gibt es nun kein separates Verzeichnis mehr für Internet- und Intranetquellen.
\item Option \verb|notkeyword=ausblenden| bei \verb|\printbibligraphy| sorgt dafür, dass Sekundärliteratur korrekt zitiert wird.
\end{itemize}

\item \verb|_dhbw_praembel.tex|
\begin{itemize}
\item (Fix) Die Bezeichnung geschachtelter Anhänge wurde auf das in den Zitierrichtlinien geforderte Format \enquote{Anhang 2/1} angepasst (Befehl \verb|\anhangteil|).
\end{itemize}

\item \verb|einleitung.tex|
\begin{itemize}
\item Hinweis zum Ausblenden der farbigen Links im PDF hinzugefügt
\end{itemize}

\item \verb|text-mit-zitaten.tex|
\begin{itemize}
\item Abschnitt 3.4 aktualisiert nach Wegfall des separaten Verzeichnisses für Internet- und Intranetquellen
\item Abschnitt zum Zitieren von Sekundärliteratur hinzugefügt
\end{itemize}

\end{itemize}

%%% Ende des eigentlichen Inhalts %%%


%%% Quellenverzeichnisse (keine Anpassungen nötig) %%%
\clearpage
\chapter*{List of references}\label{chapter:quellen}
\addcontentsline{toc}{chapter}{List of references}
\spezialkopfzeile{List of references}

% alle Quellen außer Gespräche:
\printbibliography[heading=lit,nottype=online,notkeyword=gespraech,notkeyword=ausblenden]

% alle Gespräche (diese müssen vom Typ "misc" sein und unter keyword "gespraech" stehen haben)
\printbibliography[heading=gespraeche,type=misc,keyword=gespraech]

%%% Ehrenwörtliche Erklärung (keine Anpassungen nötig) %%%
% steht ganz am Ende des Dokuments
\cleardoublepage
\clearpage

\thispagestyle{kapitelkopfzeile}

\chapter*{Erklärung}
\addcontentsline{toc}{chapter}{Ehrenwörtliche Erklärung}

% \typMeinerArbeit und \themaMeinerArbeit werden in deckblatt.tex definiert
Ich versichere hiermit, dass ich meine \typMeinerArbeit\ mit dem Thema: \emph{\themaMeinerArbeit} selbstständig verfasst und keine anderen als die angegebenen Quellen und Hilfsmittel benutzt habe. 

Ich versichere zudem, dass die eingereichte elektronische Fassung mit der gedruckten Fassung übereinstimmt.

\vspace{3cm}

\begin{center}
\begin{tabular}{ccc}
(Ort, Datum) & \hspace{0.3\linewidth} & (Unterschrift)
\end{tabular}
\end{center}

\end{document}